\section{Independent integral review of the primary cyclic lattice}
\label{sec:pclr}

The input is the specified coefficient ring of the single-primary
calculation, with its given embedding in \(\mathbb C\) and its fixed
specialization \(\alpha:R\to k_0\), where \(k_0\) is a finite field of
characteristic \(p>m+1\). The ring retains \(\rho\), \(1/(m+1)!\), every
coefficient \(a_j\) of the full Taylor unit, and every coefficient
\(b_j\) of its inverse. These satisfy the actual truncated identity
\[
 \left(\sum_{j=0}^{m-1}a_jy^j\right)
 \left(\sum_{j=0}^{m-1}b_jy^j\right)=1\quad\text{in }R[y]/(y^m).
 \tag{PCLR.1}
\]
We calculate the given specialization; existence of additional
specializations is not asserted.

Put \(\mathfrak q=\ker\alpha\), \(A=R_{\mathfrak q}\), and
\(\kappa=\alpha(R)\subseteq k_0\). The finite subring \(\alpha(R)\)
is a field: multiplication by a nonzero element is injective on its
finite underlying set and hence surjective, so that element has an
inverse. Thus \(\mathfrak q\) is maximal and
\[
 A\subseteq\mathbb C,\qquad
 A/\mathfrak qA=\kappa,\qquad
 \alpha_A(r/s)=\alpha(r)/\alpha(s).
 \tag{PCLR.2}
\]
The map is well-defined because \(s\notin\mathfrak q\). Every integer
prime to \(p\) lies outside \(\mathfrak q\) and becomes a unit in
\(A\). The integer \(p\) is a nonzero nonunit, since \(A\) is a
characteristic-zero domain and \(\alpha_A(p)=0\). No assertion that
\(A\) is a discrete valuation ring is needed.

Retain an integer \(k\geq1\), put \(K=k(m-1)\), and write
\[
 E_A=A[y_1,\ldots,y_k]/(y_1^m,\ldots,y_k^m),\qquad
 Jf=(y_1+\cdots+y_k)f.
 \tag{PCLR.3}
\]
The ordered monomials \(y^{\boldsymbol\nu}\), with
\(0\leq\nu_i\leq m-1\), form a free \(A\)-basis. In the original
coordinates \(y_i=s_i-\rho\), the Pascal matrix is
\(P_{ab}=\binom ba\rho^{b-a}\) for \(a\leq b\), zero otherwise,
and its inverse has \((-\rho)^{b-a}\) in place of \(\rho^{b-a}\).
Substitution of the two binomial expansions proves both inverse
identities over \(A\). With \(P_k=P^{\otimes k}\), the original
nilpotent is \(J_s=P_k^{-1}JP_k\), and the original sum-coordinate
action is \(k\rho I+J_s\). This retains the specified coordinates.

For \(0\leq r\leq K\), define the following element of the original
marked algebra, using only invertible factorials:
\[
 V_r=\sum_{\substack{0\leq\nu_i<m\\|\boldsymbol\nu|=r}}
          \frac{y^{\boldsymbol\nu}}{\nu_1!\cdots\nu_k!},
 \qquad W=\bigoplus_{r=0}^K A V_r.
 \tag{PCLR.4}
\]
Every index \(\nu_i\leq m-1<p\), so the displayed denominators are
units. For each \(r\) choose \(\boldsymbol\nu(r)\) greedily: starting
with remainder \(r\), set \(\nu_1=\min(m-1,r)\), subtract it, and
continue. Since \(r\leq K\), this produces a permitted tuple of
total degree \(r\). When \(m=1\), the only case is \(r=0\) and the
tuple is zero. If \([y^{\boldsymbol\nu}]f\) denotes the coefficient
of a basis monomial, put
\[
 \lambda_r(f)=\left(\prod_i\nu_i(r)!\right)
                    [y^{\boldsymbol\nu(r)}]f,\qquad
 \pi(f)=\sum_{r=0}^K\lambda_r(f)V_r.
 \tag{PCLR.5}
\]
Different \(V_r\) have different degrees, and the selected
coefficient of \(V_r\) is the inverse of the product in
\(\lambda_r\). Consequently \(\lambda_r(V_t)=\delta_{rt}\),
\(\pi^2=\pi\), and
\[
 E_A=W\oplus\ker\pi.
 \tag{PCLR.6}
\]
More explicitly, \(\ker\pi\) is freely spanned by all original
monomials other than the \(K+1\) selected monomials: the condition
\(\pi(f)=0\) is exactly the vanishing of those selected
coefficients. Thus \(W\) is a primitive split summand, including
at the specified characteristic-\(p\) fibre.

The multinomial expansion and multiplication of (PCLR.4) give
\[
 J^r1=r!V_r\quad(0\leq r\leq K),\qquad
 JV_r=(r+1)V_{r+1}\quad(0\leq r<K),\qquad JV_K=0.
 \tag{PCLR.7}
\]
For the middle identity, the coefficient of a retained monomial
\(y^{\boldsymbol\mu}\) of degree \(r+1\) is
\(\sum_{i:\mu_i>0}1/(\mu_i-1)!\prod_{j\ne i}1/\mu_j!
=(\sum_i\mu_i)/\prod_i\mu_i!\).
Terms outside the retained exponent range vanish in the quotient.
There are no retained monomials of degree greater than \(K\),
which proves the last identity.

The ordinary cyclic lattice, as an actual submodule of \(E_A\), is
\[
 L=A[J]1=\bigoplus_{r=0}^K A\ell_r,\qquad
 \ell_r=r!V_r,\qquad
 J\ell_r=\ell_{r+1}\ (r<K),\quad J\ell_K=0.
 \tag{PCLR.8}
\]
Independence follows from (PCLR.6) and the fact that each \(r!\)
is a nonzero element of the domain \(A\). In these displayed bases
the inclusion \(L\hookrightarrow W\) is diagonal, with entries
\(0!,1!,\ldots,K!\). Therefore
\[
 T=W/L=\bigoplus_{r=0}^K A/(r!),\qquad
 E_A/L\simeq\ker\pi\oplus T.
 \tag{PCLR.9}
\]
The equality denotes the explicit map taking
\(\sum c_rV_r\) to \((c_r\bmod r!)_r\), whose kernel is (PCLR.8).
Writing \(e_r=\sum_{j\geq1}\lfloor r/p^j\rfloor\), counting the
factors divisible by each \(p^j\) in \(r!\) gives
\(r!=p^{e_r}\varepsilon_r\), with \(\varepsilon_r\) an integer
prime to \(p\). As \(\varepsilon_r\in A^\times\),
\[
 A/(r!)=A/(p^{e_r}),\qquad
 T=\bigoplus_{r=p}^K A/(p^{e_r}).
 \tag{PCLR.10}
\]
An empty direct sum is zero. This computes torsion over the actual
localized coefficient ring without replacing that ring by
\(\mathbb Z_{(p)}\).

Let \(U\) denote multiplication by
\[
 \prod_{i=1}^k\left(\sum_{j=0}^{m-1}a_jy_i^j\right).
 \tag{PCLR.11}
\]
Multiplication by the product of the corresponding \(b_j\) series
is its inverse by (PCLR.1), and \(UJ=JU\) because \(E_A\) is
commutative. Put \(W_U=UW\), \(L_U=UL\), and
\(\pi_U=U\pi U^{-1}\). Then the identities just proved give
\[
 E_A=W_U\oplus U\ker\pi,\qquad
 L_U=A[J](U1),\qquad
 W_U/L_U\xrightarrow[\ f\mapsto U^{-1}f\ ]{\sim}W/L.
 \tag{PCLR.12}
\]
The inclusion is still the diagonal matrix \(r!\) in the exact
bases \(U\ell_r\) and \(UV_r\). This transports the complete unit;
it does not assert that \(U\) preserves \(W\) or \(L\).

The base-change calculation is a calculation of a specific free
resolution. Write \(\overline M=M\otimes_A\kappa\), and put
\(d=\min(p,K+1)\). The resolution
\(0\to L\to W\to T\to0\), in the bases above, becomes the two-term
complex with matrix \(\operatorname{diag}(\overline{r!})\).
Consequently its homology gives the exact sequence and named maps
\[
 0\longrightarrow\operatorname{Tor}_1^A(T,\kappa)
 \xrightarrow{\delta}\overline L
 \xrightarrow{\iota_\kappa}\overline W
 \longrightarrow T\otimes_A\kappa\longrightarrow0,
 \tag{PCLR.13}
\]
\[
 \begin{aligned}
 \operatorname{Tor}_1^A(T,\kappa)
    &\simeq\bigoplus_{r=p}^K\kappa,
 &\delta((z_r)_r)&=\sum_{r=p}^Kz_r\overline\ell_r,\\
 \iota_\kappa(\overline\ell_r)&=
                  \overline{r!}\,\overline V_r,
 &\operatorname{im}\iota_\kappa&=
                  \bigoplus_{r=0}^{d-1}\kappa\overline V_r,\\
 T\otimes_A\kappa&\simeq\bigoplus_{r=p}^K\kappa,
 &\dim_\kappa\operatorname{im}\iota_\kappa&=d.
 \end{aligned}
 \tag{PCLR.14}
\]
The displayed Tor identification uses the resolution with
multiplication by \(r!\); replacing it by a resolution with
\(p^{e_r}\) rescales its degree-one identification by a unit.
For completeness, \(r!\) is invertible in \(\kappa\) exactly when
\(r<p\), and is zero otherwise. Thus each diagonal coordinate
has zero kernel and cokernel for \(r<p\), and one-dimensional
kernel and cokernel for \(r\geq p\). This proves all maps in
(PCLR.13)--(PCLR.14), rather than only their dimensions. Extension
of scalars from \(\kappa\) to the originally specified \(k_0\)
is exact: choosing a \(\kappa\)-basis of \(k_0\) identifies
tensoring with a direct sum of copies of the original vector
space. The same formulas therefore hold over \(k_0\).

There are two different endomorphisms after this base change,
related by the map in (PCLR.13). On \(\overline E_A\), the binomial
formula in characteristic \(p\) gives
\[
 \overline J^{\,p}=\sum_i y_i^p=0,
 \qquad\overline J^{\,K+1}=0,
 \qquad \overline J^{\,d-1}1=(d-1)!\,\overline V_{d-1}\ne0.
 \tag{PCLR.15}
\]
The first equality follows because all intermediate binomial
coefficients are divisible by the prime \(p\), iterated over
the commuting variables. Each \(y_i^p\) is zero since \(p>m\).
The nonvanishing uses the split summand and \(d-1<p\). Hence the
ambient endomorphism and its restriction to the cyclic image
both have exact nilpotency index \(d\). Its restriction to
\(\overline W\) has the same index, and decomposes into the chains
\[
 \operatorname{span}_\kappa
  \{\overline V_{jp},\ldots,
     \overline V_{\min((j+1)p-1,K)}\},
 \qquad 0\leq j\leq\lfloor K/p\rfloor.
 \tag{PCLR.16}
\]
Inside each displayed chain, the coefficients \(r+1\) in
(PCLR.7) are nonzero until its endpoint, and the next coefficient
is divisible by \(p\), unless the last endpoint is \(K\), where
the action is zero by degree. These are all the basis vectors,
so the chain decomposition is exact.

In contrast, the endomorphism obtained by first restricting
\(J\) to \(L\), then tensoring, satisfies
\[
 J_{\overline L}\overline\ell_r=\overline\ell_{r+1},\qquad
 (\overline L,J_{\overline L})\simeq
           (\kappa[t]/(t^{K+1}),\text{multiplication by }t).
 \tag{PCLR.17}
\]
Its nilpotency index is \(K+1\), including when \(K\geq p\).
The exact intertwiner \(\iota_\kappa\) sends the last \(K+1-d\)
chain vectors to zero. In polynomial notation its kernel and
image are precisely
\[
 \ker\iota_\kappa=t^d\kappa[t]/(t^{K+1}),\qquad
 \operatorname{im}\iota_\kappa\simeq\kappa[t]/(t^d).
 \tag{PCLR.18}
\]
For \(K\geq p\) this also proves directly that the kernel is a
single nilpotent chain of length \(K+1-p\). There is no
contradiction with (PCLR.15): the inclusion ceased to be
injective after tensoring. All identities (PCLR.13)--(PCLR.18)
transport by \(\overline U\), an invertible matrix because its
specified inverse also specializes, and by the retained Pascal
matrices.

The terminal monomial and its exact scalar are
\[
 \tau_k=y_1^{m-1}\cdots y_k^{m-1},\quad
 V_K=\frac{\tau_k}{((m-1)!)^k},\quad
 U\ell_K=\frac{K!}{((m-1)!)^k}a_0^k\tau_k.
 \tag{PCLR.19}
\]
Only the tuple \((m-1,\ldots,m-1)\) occurs in degree \(K\).
Every positive-degree term of the unit kills that monomial;
its surviving factor is exactly \(a_0^k\), with
\(a_0b_0=1\). The class \(\overline\tau_k\) is a nonzero original
basis vector, and \(\overline{a_0}\ne0\). Hence the displayed
specialized terminal scalar vanishes exactly when \(K\geq p\).
When \(m=1\), \(K=0\), \(L=W=E_A\), all Tor terms vanish and
the scalar is \(a_0^k\); every formula above includes this case.

Finally the marked source already supplies the fixed-nonbottom
split lift \(\mathsf S(M)=\{\tau\}\sqcup M^\bullet\).
For every actual coefficient map \(f\) calculated here its lift
is exactly
\[
 \widetilde f(\tau)=\tau,\qquad
 \widetilde f(x^\bullet)=f(x)^\bullet,\qquad
 \widetilde f^{-1}(0^\bullet)=(\ker f)^\bullet,\qquad
 \widetilde f^{-1}(\tau)=\{\tau\}.
 \tag{PCLR.20}
\]
These statements follow by evaluating the two disjoint types
of elements. Composition is preserved because
\(g(f(x))=(gf)(x)\). Thus in (PCLR.14) the specialized high
chain vectors map to the represented zero at the same support
label, and never to external absence. The original marked note
specifies this lift at fixed nonbottom labels; no extra
support-changing transport map is introduced by this review.

\subsection{The exact higher-support chain supplied by the factorial matrix}

Put \(D=e_K=v_p(K!)\). When \(D=0\), the quotient \(T\) in
(PCLR.9) is zero and no positive-depth chain is required. For
\(D>0\), put \(A_D=A/p^DA\) and
\(B_D=\mathbb Z/p^D\mathbb Z\). The integer structural map gives
\[
 \iota_D:B_D\hookrightarrow A_D.
 \tag{PCLR.21}
\]
Here is a proof of injectivity over this actual ring. Suppose a
nonzero residue integer \(n\bmod p^D\) lies in its kernel. Factor
\(n=p^vu\), where \(0\leq v<D\) and \(p\nmid u\). Membership
\(n\in p^DA\) and cancellation of the nonzero \(p^v\) in the
domain \(A\) give \(u=p^{D-v}b\) for some \(b\in A\). The left
side is a unit, while the right side lies in the maximal ideal
\(\mathfrak qA\). This is impossible, proving (PCLR.21).

For \(0\leq a\leq D\), set \(I_a=p^aB_D\) and
\(\mathfrak P_a=p^aA_D\). The actual extension and contraction
maps satisfy
\[
 \iota_D(I_a)A_D=\mathfrak P_a,
 \qquad \iota_D^{-1}(\mathfrak P_a)=I_a,
 \qquad \mathfrak P_a\supsetneqq\mathfrak P_{a+1}
                 \quad(a<D).
 \tag{PCLR.22}
\]
Extension follows from the displayed generators. For contraction,
an integer with valuation \(v<a\) in \(p^aA\) gives the same
unit-in-the-maximal-ideal contradiction after cancellation;
integers divisible by \(p^a\) are in the contraction directly.
For strictness, an equality
\(p^a=p^{a+1}b\bmod p^DA\) would give
\(1=pb+p^{D-a}c\) in \(A\) after cancellation. Its right side
lies in \(\mathfrak qA\), a contradiction. This also shows
\(A_D\ne0\). Because the displayed ideals are nested, and
products of principal ideals have the product generator,
\[
 \begin{split}
 \mathfrak P_a+\mathfrak P_b&=\mathfrak P_{\min(a,b)},\qquad
 \mathfrak P_a\cap\mathfrak P_b=\mathfrak P_{\max(a,b)},\\
 \mathfrak P_a\mathfrak P_b&=\mathfrak P_{\min(D,a+b)}.
 \end{split}
 \tag{PCLR.23}
\]
Thus extension preserves sums, intersections, and ideal products
on this precise subchain. We do not assert that every ideal of
\(A_D\) is one of these principal ideals.

Use now the original MFC.29 carrier for the original chain
\(C_D=\{0<\cdots<D\}\):
\[
 \begin{split}
 G_{C_D}(A_D)&=\{(0,j):j\in C_D\}
                      \cup\{(a,D):a\in A_D\},\\
 (a,j)+(b,l)&=(a+b,\max(j,l)),\\
 (a,j)(b,l)&=(ab,\min(j,l)).
 \end{split}
 \tag{PCLR.24}
\]
The structural coefficient injection in (PCLR.21) has the exact
support-preserving lift
\[
 G_{C_D}(B_D)\hookrightarrow G_{C_D}(A_D),
       \qquad (b,j)\longmapsto(\iota_D(b),j).
 \tag{PCLR.25}
\]
The two permitted kinds of pairs map to the same respective
kinds. Substitution in the addition and multiplication formulas
proves both laws, and the map preserves \((0,0)\) and \((1,D)\).
Injectivity follows from injectivity of \(\iota_D\) and equality
of the support coordinate.

Let \(\mathcal P_D(A)=\{\mathfrak P_0,\ldots,
\mathfrak P_D\}\), with addition given by ideal sum and
multiplication by ideal intersection. Its zero is
\(\mathfrak P_D\) and its unit is \(\mathfrak P_0\). The original
support projection followed by the exact ideal-chain map is
\[
 G_{C_D}(A_D)\xrightarrow{\chi}(C_D,\max,\min)
    \xrightarrow{\theta_A}\mathcal P_D(A),\qquad
 \chi(a,j)=j,\quad\theta_A(j)=\mathfrak P_{D-j}.
 \tag{PCLR.26}
\]
Equation (PCLR.24) proves that \(\chi\) preserves both
operations. Strictness and the sum/intersection laws in
(PCLR.22)--(PCLR.23) prove that \(\theta_A\) is a bijective
homomorphism of bounded-lattice semirings. In the original
MFC.35 indexing \(\tau_a=(0,D-a)\), including
\(\tau_0=e=(0,D)\), this composite is exactly
\[
 \theta_A\chi(\tau_a)=\mathfrak P_a\quad(0\leq a\leq D).
 \tag{PCLR.27}
\]
The original ideal-product operation is also retained: its
transport to support coordinates is
\(j\mathbin{\odot}l=\max(0,j+l-D)\), because
\(D-\min(D,(D-j)+(D-l))=\max(0,j+l-D)\).
This operation is different from the \(\min\) in the carrier
multiplication; (PCLR.26) uses intersection, as its declared
codomain specifies.

Finally this support chain measures the original cyclic map by
an exact coordinate morphism. Define
\[
 f_r:A[t]/(t^{K+1})\longrightarrow A,\qquad
 f_r(F)=\lambda_rU^{-1}\bigl(F(J)U1\bigr).
 \tag{PCLR.28}
\]
The map is \(A\)-linear, since every constituent is, and
\(U^{-1}F(J)U1=F(J)1\). Equations (PCLR.5) and (PCLR.7)
therefore give \(f_r(t^s)=\delta_{rs}r!\) for all
\(0\leq s\leq K\). It follows that
\[
 \operatorname{im}f_r=(r!)=p^{e_r}A,\qquad
 \operatorname{im}(f_r\otimes_A A_D)=\mathfrak P_{e_r}.
 \tag{PCLR.29}
\]
The one-generator cokernel has presentation by the matrix
\([r!]\), so the ideal generated by its maximal minors, its
zeroth Fitting ideal, is exactly \((r!)\). The label of this
coordinate ideal under (PCLR.26)--(PCLR.27) is therefore the
specified original stratum \(\tau_{e_r}\). This is a proved
coordinate observation of the actual full-unit cyclic map;
it is not an invented transport sending arbitrary coefficient
elements to support strata.

\paragraph{Input pins.}
This proof uses the retained coefficient datum and complete unit
of SPF.1, the ordered original algebra and Pascal matrix TP.1--2,
the original full unit TP.23, and the fixed-nonbottom split lift
of the original marked note, equation (8), together with the
original higher-support carrier and chain maps MFC.29--44.
Every algebraic
specialization, inclusion, retraction, quotient, Tor map and
nilpotency statement used above is proved explicitly here.
